% !TeX program = lualatex
\documentclass[12pt,aspectratio=169]{beamer}
\usepackage[utf8]{inputenc}
\usepackage[T1]{fontenc}
\usepackage{amsmath,amsfonts,amssymb}
\usepackage{graphicx}
\usepackage{xcolor}
\usepackage{hyperref}
\usepackage{anyfontsize}
\usepackage{lmodern}

% ====== EXTREME FONT REDUCTION ======
\renewcommand{\tiny}{\fontsize{4}{5}\selectfont}
\renewcommand{\scriptsize}{\fontsize{5}{6}\selectfont}
\renewcommand{\footnotesize}{\fontsize{6}{7}\selectfont}
\renewcommand{\small}{\fontsize{7}{8}\selectfont}
\renewcommand{\normalsize}{\fontsize{8}{9}\selectfont}
\renewcommand{\large}{\fontsize{9}{10}\selectfont}
\renewcommand{\Large}{\fontsize{10}{11}\selectfont}
\renewcommand{\LARGE}{\fontsize{11}{13}\selectfont}
\renewcommand{\huge}{\fontsize{12}{14}\selectfont}
\renewcommand{\Huge}{\fontsize{14}{17}\selectfont}

% ====== COLOR THEME ======
\definecolor{redcorrect}{RGB}{200,30,30}
\definecolor{bluecorrect}{RGB}{30,80,200}
\definecolor{greencheck}{RGB}{0,150,50}
\definecolor{lightgray}{RGB}{245,245,245}
\definecolor{darkgray}{RGB}{50,50,50}
\definecolor{softyellow}{RGB}{255,248,200}
\definecolor{steppurple}{RGB}{150,50,200}
\definecolor{deepblue}{RGB}{20,60,150}
\definecolor{darkred}{RGB}{150,20,20}
\definecolor{orangewarning}{RGB}{255,140,0}

\setbeamercolor{title}{fg=white,bg=redcorrect}
\setbeamercolor{frametitle}{fg=white,bg=redcorrect}
\setbeamercolor{block title}{fg=white,bg=bluecorrect}
\setbeamercolor{block body}{fg=darkgray,bg=lightgray}
\setbeamercolor{normal text}{fg=darkgray}
\usecolortheme{default}

% ====== CUSTOM COMMANDS ======
\newcommand{\correct}[1]{\textcolor{greencheck}{\textbf{\checkmark}} #1}
\newcommand{\keypoint}[1]{\textcolor{deepblue}{\textbf{#1}}}
\newcommand{\hardconcept}[1]{\textcolor{darkred}{\textbf{#1}}}
\newcommand{\tipbox}[1]{\fcolorbox{bluecorrect}{softyellow}{\parbox{0.88\textwidth}{\centering #1}}}
\newcommand{\stepbox}[1]{%
	\begin{center}
		\fcolorbox{darkgray}{white}{%
			\parbox{0.90\textwidth}{\vspace{0.05cm} #1 \vspace{0.05cm}}%
		}
	\end{center}
}
\newcommand{\wrongbox}[1]{\fcolorbox{redcorrect}{red!10}{\parbox{0.88\textwidth}{\centering \textcolor{redcorrect}{\textbf{\times}} #1}}}
\newcommand{\highlight}[1]{\textcolor{redcorrect}{\textbf{#1}}}
\newcommand{\bluetext}[1]{\textcolor{bluecorrect}{\textbf{#1}}}

\begin{document}
	
	% ================================================================
	% SLIDE 1-2: TITLE AND AGENDA
	% ================================================================
	\begin{frame}
		\centering
		\vspace{0.8cm}
		{\fontsize{16}{19}\selectfont \textbf{Global AFC Frameworks, Governance, and Regulations}}\\[0.15cm]
		{\fontsize{12}{14}\selectfont Domain B — Complete Q\&A Review}\\[0.3cm]
		{\fontsize{9}{11}\selectfont Understanding the \textit{Regulatory Principles} Behind Each Answer}\\[0.3cm]
		{\fontsize{8}{10}\selectfont Compliance Training Department}
		\vspace{0.2cm}
		\rule{4cm}{0.5pt}
	\end{frame}
	
	\begin{frame}
		\frametitle{\fontsize{11}{13}\selectfont Agenda — Domain B Topics}
		\scriptsize
		\begin{enumerate}
			\item FATF Role, Function, and Regional Bodies (Q1-Q6)
			\item UN Role and Sanctions Regimes (Q7-Q10)
			\item Regulators, Law Enforcement, and FIUs (Q11-Q16)
			\item Public Sector vs Private Sector Objectives (Q17-Q20)
			\item Regulatory Cooperation — Cross-Border (Q21-Q24)
			\item Leveraging Regulatory Reports (Q25-Q28)
			\item Non-Government Body Guidance (Q29-Q32)
			\item National and Sectoral Risk Assessments (Q33-Q36)
			\item US and EU AFC Regulations (Q37-Q42)
			\item Jurisdictional Regulatory Awareness (Q43-Q46)
			\item Public-Private Partnerships — Why Important (Q47-Q50)
			\item PPP — Strengths and Limitations (Q51-Q54)
			\item Private Sector Collaboration (Q55-Q58)
			\item Data and Intelligence Sharing Considerations (Q59-Q62)
			\item Other Laws Impacting AML — Data Privacy (Q63-Q66)
			\item Other Laws — Consumer Protection (Q67-Q70)
			\item Other Laws — ESG (Q71-Q74)
			\item Other Laws — Financial Inclusion (Q75-Q78)
		\end{enumerate}
	\end{frame}
	
	% ================================================================
	% SLIDE 3-8: FATF ROLE AND FUNCTION
	% ================================================================
	\begin{frame}
		\frametitle{\fontsize{11}{13}\selectfont Q1: FATF Role and Function}
		\begin{block}{\fontsize{8}{10}\selectfont Topic: FATF Role, Function, and Regional Bodies}
			\scriptsize
			Which of the following best describes the primary function of the Financial Action Task Force (FATF)?
		\end{block}
		\vspace{0.05cm}
		\stepbox{\scriptsize
			\keypoint{Correct Answer:} Setting global standards for AML/CFT and conducting mutual evaluations to assess member compliance.
		}
		\vspace{0.05cm}
		\scriptsize
		\hardconcept{Core Concept:} The FATF is an \textbf{intergovernmental body} established in 1989. Its primary functions are:
		\par \textbf{Standard Setting}: Issuing the 40 Recommendations that form the global AML/CFT framework.
		\par \textbf{Mutual Evaluations}: Assessing member countries' compliance through peer reviews.
		\par \textbf{Typology Research}: Identifying and publishing new money laundering and terrorist financing methods.
		\par \textbf{Black/Grey Listing}: Identifying jurisdictions with strategic AML/CFT deficiencies.
		\vspace{0.03cm}
		\wrongbox{\scriptsize \textbf{Common Trap:} "FATF enforces AML laws directly" — INCORRECT. FATF sets standards but does not enforce laws; enforcement is done by national regulators.}
		\vspace{0.03cm}
		\tipbox{\scriptsize \textbf{Key Insight:} FATF Recommendations are not legally binding but are implemented through national legislation. Countries that fail to comply risk being placed on the FATF Grey List.}
	\end{frame}
	
	\begin{frame}
		\frametitle{\fontsize{11}{13}\selectfont Q2: FATF's 40 Recommendations}
		\begin{block}{\fontsize{8}{10}\selectfont Topic: FATF 40 Recommendations}
			\scriptsize
			How many Recommendations does the FATF issue, and what do they cover?
		\end{block}
		\vspace{0.05cm}
		\stepbox{\scriptsize
			\keypoint{Correct Answer:} 40 Recommendations covering AML/CFT, including customer due diligence, record-keeping, suspicious transaction reporting, sanctions, and international cooperation.
		}
		\vspace{0.05cm}
		\scriptsize
		\hardconcept{Core Concept:} The 40 Recommendations are organized into:
		\par \textbf{Recommendations 1-2}: AML/CFT policies and coordination.
		\par \textbf{Recommendations 3-8}: Money laundering and terrorist financing offenses.
		\par \textbf{Recommendations 9-23}: Preventive measures for financial institutions and DNFBPs.
		\par \textbf{Recommendations 24-25}: Transparency of legal persons and arrangements.
		\par \textbf{Recommendations 26-35}: Powers and responsibilities of competent authorities.
		\par \textbf{Recommendations 36-40}: International cooperation.
		\vspace{0.03cm}
		\tipbox{\scriptsize \textbf{Key Insight:} \textbf{Recommendation 1} establishes the risk-based approach. \textbf{Recommendation 10} covers Customer Due Diligence. \textbf{Recommendation 20} covers STR filing.}
	\end{frame}
	
	\begin{frame}
		\frametitle{\fontsize{11}{13}\selectfont Q3: FATF Mutual Evaluations}
		\begin{block}{\fontsize{8}{10}\selectfont Topic: FATF Mutual Evaluations}
			\scriptsize
			What do FATF mutual evaluations assess?
		\end{block}
		\vspace{0.05cm}
		\stepbox{\scriptsize
			\keypoint{Correct Answer:} Both technical compliance (whether laws exist) and effectiveness (whether laws produce results against 11 Immediate Outcomes).
		}
		\vspace{0.05cm}
		\scriptsize
		\hardconcept{Core Concept:}
		\par \textbf{Technical Compliance}: Does the country have the required laws and regulations?
		\par \textbf{Effectiveness}: Are the laws producing results? Measured against \textbf{11 Immediate Outcomes}.
		\par \textbf{IO.1}: ML/TF risks understood.
		\par \textbf{IO.4}: AML/CFT controls implemented.
		\par \textbf{IO.7}: ML investigated and prosecuted.
		\vspace{0.03cm}
		\tipbox{\scriptsize \textbf{Key Insight:} Effectiveness ratings are more important than technical compliance ratings. A country can have perfect laws (technical compliance) but fail to implement them (low effectiveness).}
	\end{frame}
	
	\begin{frame}
		\frametitle{\fontsize{11}{13}\selectfont Q4: FATF Grey and Black Lists}
		\begin{block}{\fontsize{8}{10}\selectfont Topic: FATF Grey and Black Lists}
			\scriptsize
			What is the difference between FATF's Grey List and Black List?
		\end{block}
		\vspace{0.05cm}
		\stepbox{\scriptsize
			\keypoint{Correct Answer:} The Black List identifies jurisdictions with strategic AML/CFT deficiencies that have not made sufficient progress; the Grey List identifies jurisdictions under increased monitoring that are actively working to address deficiencies.
		}
		\vspace{0.05cm}
		\scriptsize
		\hardconcept{Core Concept:}
		\par \textbf{Black List (Call for Action)}: Jurisdictions with serious strategic deficiencies. Countermeasures are recommended.
		\par \textbf{Grey List (Increased Monitoring)}: Jurisdictions with strategic deficiencies that have committed to an action plan. They are actively working to address issues.
		\vspace{0.03cm}
		\tipbox{\scriptsize \textbf{Key Insight:} Being on the Grey List has significant reputational and economic consequences. Financial institutions must apply enhanced due diligence to clients from Grey List jurisdictions.}
	\end{frame}
	
	\begin{frame}
		\frametitle{\fontsize{11}{13}\selectfont Q5: FATF-Style Regional Bodies (FSRBs)}
		\begin{block}{\fontsize{8}{10}\selectfont Topic: FATF-Style Regional Bodies}
			\scriptsize
			Which of the following is a FATF-Style Regional Body (FSRB)?
		\end{block}
		\vspace{0.05cm}
		\stepbox{\scriptsize
			\keypoint{Correct Answer:} APG (Asia/Pacific Group on Money Laundering), GAFILAT (Latin America), MONEYVAL (Europe), ESAAMLG (Eastern and Southern Africa), and MENAFATF (Middle East and North Africa).
		}
		\vspace{0.05cm}
		\scriptsize
		\hardconcept{Core Concept:} FSRBs are regional organizations that:
		\par \textbf{Implement FATF standards}: Adapt global standards to regional contexts.
		\par \textbf{Conduct mutual evaluations}: Assess member compliance.
		\par \textbf{Provide technical assistance}: Help countries implement frameworks.
		\par \textbf{Promote cooperation}: Facilitate regional information sharing.
		\vspace{0.03cm}
		\wrongbox{\scriptsize \textbf{Common Trap:} "FSRBs are independent from FATF and set their own standards" — INCORRECT. FSRBs are \textbf{affiliated with} and \textbf{implement} FATF standards.}
		\vspace{0.03cm}
		\tipbox{\scriptsize \textbf{Key Insight:} There are nine FSRBs worldwide. Each is a member of the FATF Global Network and participates in the FATF Plenary.}
	\end{frame}
	
	\begin{frame}
		\frametitle{\fontsize{11}{13}\selectfont Q6: Applying FATF Recommendations}
		\begin{block}{\fontsize{8}{10}\selectfont Topic: Applying FATF Recommendations}
			\scriptsize
			A financial institution is updating its AML program. Which approach best reflects the proper application of FATF Recommendations?
		\end{block}
		\vspace{0.05cm}
		\stepbox{\scriptsize
			\keypoint{Correct Answer:} Implement a risk-based approach that tailors controls to the institution's specific risk profile while meeting the core requirements of each Recommendation.
		}
		\vspace{0.05cm}
		\scriptsize
		\hardconcept{Core Concept:} The \textbf{risk-based approach} is the foundation of FATF Recommendations. It requires:
		\par \textbf{Risk Assessment}: Identify and assess ML/TF risks.
		\par \textbf{Proportionate Controls}: Implement controls appropriate to the level of risk.
		\par \textbf{Continuous Monitoring}: Regularly review and update the program.
		\vspace{0.03cm}
		\wrongbox{\scriptsize \textbf{Common Trap:} "Apply all FATF Recommendations uniformly without deviation" — INCORRECT. FATF explicitly requires a \textbf{risk-based approach}, not a one-size-fits-all approach.}
		\vspace{0.03cm}
		\tipbox{\scriptsize \textbf{Key Insight:} \textbf{FATF Recommendation 1} establishes the risk-based approach. It requires countries and institutions to identify, assess, and understand their ML/TF risks and apply resources proportionately.}
	\end{frame}
	
	% ================================================================
	% SLIDE 9-12: UN ROLE AND SANCTIONS
	% ================================================================
	\begin{frame}
		\frametitle{\fontsize{11}{13}\selectfont Q7: UN Role in AML/CFT}
		\begin{block}{\fontsize{8}{10}\selectfont Topic: UN Role and Sanctions Regimes}
			\scriptsize
			Which of the following accurately describes the role of the United Nations in AML/CFT?
		\end{block}
		\vspace{0.05cm}
		\stepbox{\scriptsize
			\keypoint{Correct Answer:} The UN establishes binding sanctions regimes through Security Council resolutions and promotes AML/CFT through conventions like the UN Convention against Corruption (UNCAC).
		}
		\vspace{0.05cm}
		\scriptsize
		\hardconcept{Core Concept:} The UN plays several key roles:
		\par \textbf{Sanctions Regimes}: UN Security Council resolutions impose binding sanctions.
		\par \textbf{Conventions}: UNTOC (Transnational Organized Crime) and UNCAC (Corruption) provide legal frameworks.
		\par \textbf{Coordination}: UNODC provides technical assistance.
		\vspace{0.03cm}
		\wrongbox{\scriptsize \textbf{Common Trap:} "UN sanctions are optional recommendations" — INCORRECT. UN Security Council resolutions are \textbf{binding} on all member states under Article 25 of the UN Charter.}
		\vspace{0.03cm}
		\tipbox{\scriptsize \textbf{Key Insight:} \textbf{UNSCR 1267} (Al-Qaida/Taliban) and \textbf{UNSCR 1373} (terrorist financing) are the foundation of UN sanctions regimes.}
	\end{frame}
	
	\begin{frame}
		\frametitle{\fontsize{11}{13}\selectfont Q8: UNSCR 1267 and 1373}
		\begin{block}{\fontsize{8}{10}\selectfont Topic: UN Sanctions Resolutions}
			\scriptsize
			What is the difference between UNSCR 1267 and UNSCR 1373?
		\end{block}
		\vspace{0.05cm}
		\stepbox{\scriptsize
			\keypoint{Correct Answer:} UNSCR 1267 establishes a specific sanctions list for individuals and entities associated with Al-Qaida and the Taliban; UNSCR 1373 requires all member states to criminalize terrorist financing and freeze assets of terrorists and their supporters.
		}
		\vspace{0.05cm}
		\scriptsize
		\hardconcept{Core Concept:}
		\par \textbf{UNSCR 1267 (1999)}: Created the Al-Qaida/Taliban sanctions list. Requires all states to freeze assets of listed individuals and entities.
		\par \textbf{UNSCR 1373 (2001)}: Requires all states to criminalize terrorist financing, freeze assets, and prevent the movement of terrorists.
		\vspace{0.03cm}
		\tipbox{\scriptsize \textbf{Key Insight:} UNSCR 1267 is a specific list; UNSCR 1373 requires countries to create their own domestic lists. Both require asset freezing without delay.}
	\end{frame}
	
	\begin{frame}
		\frametitle{\fontsize{11}{13}\selectfont Q9: UN Convention against Corruption (UNCAC)}
		\begin{block}{\fontsize{8}{10}\selectfont Topic: UNCAC}
			\scriptsize
			What is the primary purpose of the UN Convention against Corruption (UNCAC)?
		\end{block}
		\vspace{0.05cm}
		\stepbox{\scriptsize
			\keypoint{Correct Answer:} To provide a comprehensive global framework for preventing, detecting, and prosecuting corruption, including provisions for asset recovery and international cooperation.
		}
		\vspace{0.05cm}
		\scriptsize
		\hardconcept{Core Concept:} UNCAC (2003) covers:
		\par \textbf{Prevention}: Anti-corruption bodies, public sector transparency.
		\par \textbf{Criminalization}: Bribery, embezzlement, money laundering.
		\par \textbf{Asset Recovery}: Freezing, seizure, and return of stolen assets.
		\par \textbf{International Cooperation}: MLA, extradition, law enforcement cooperation.
		\vspace{0.03cm}
		\tipbox{\scriptsize \textbf{Key Insight:} UNCAC is the \textbf{first legally binding global anti-corruption instrument}. It has 190+ parties. Asset recovery is a core principle.}
	\end{frame}
	
	\begin{frame}
		\frametitle{\fontsize{11}{13}\selectfont Q10: UNODC Role}
		\begin{block}{\fontsize{8}{10}\selectfont Topic: UNODC}
			\scriptsize
			What is the role of the UN Office on Drugs and Crime (UNODC) in AML/CFT?
		\end{block}
		\vspace{0.05cm}
		\stepbox{\scriptsize
			\keypoint{Correct Answer:} Providing technical assistance, training, and guidance to countries implementing AML/CFT frameworks and international conventions.
		}
		\vspace{0.05cm}
		\scriptsize
		\hardconcept{Core Concept:} UNODC:
		\par \textbf{Provides Technical Assistance}: Helps countries implement AML/CFT laws.
		\par \textbf{Delivers Training}: Builds capacity in law enforcement and financial intelligence.
		\par \textbf{Publishes Guidance}: On AML/CFT best practices.
		\par \textbf{Monitors Implementation}: Tracks progress on UN conventions.
		\vspace{0.03cm}
		\tipbox{\scriptsize \textbf{Key Insight:} UNODC works with FIUs, law enforcement, and financial regulators to strengthen AML/CFT frameworks globally.}
	\end{frame}
	
	% ================================================================
	% SLIDE 13-18: REGULATORS, LE, FIUS
	% ================================================================
	\begin{frame}
		\frametitle{\fontsize{11}{13}\selectfont Q11: Regulators — Role}
		\begin{block}{\fontsize{8}{10}\selectfont Topic: Regulators, Law Enforcement, and FIUs}
			\scriptsize
			What is the primary role of financial regulators in AML/CFT?
		\end{block}
		\vspace{0.05cm}
		\stepbox{\scriptsize
			\keypoint{Correct Answer:} Setting AML/CFT rules, supervising compliance, and enforcing sanctions against non-compliant institutions.
		}
		\vspace{0.05cm}
		\scriptsize
		\hardconcept{Core Concept:} Regulators:
		\par \textbf{Set Rules}: Issue regulations (e.g., FinCEN, FCA, BaFin).
		\par \textbf{Supervise}: Conduct examinations and inspections.
		\par \textbf{Enforce}: Issue fines, penalties, and enforcement actions.
		\par \textbf{Guidance}: Issue guidance on expectations.
		\vspace{0.03cm}
		\wrongbox{\scriptsize \textbf{Common Trap:} "Regulators investigate and prosecute financial crimes" — INCORRECT. Regulators oversee compliance; law enforcement investigates crimes.}
		\vspace{0.03cm}
		\tipbox{\scriptsize \textbf{Key Insight:} Key regulators include FinCEN (US), FCA (UK), BaFin (Germany), AMF (France), and the Fed (US banking).}
	\end{frame}
	
	\begin{frame}
		\frametitle{\fontsize{11}{13}\selectfont Q12: Financial Intelligence Units (FIUs)}
		\begin{block}{\fontsize{8}{10}\selectfont Topic: Financial Intelligence Units}
			\scriptsize
			What is the primary function of a Financial Intelligence Unit (FIU)?
		\end{block}
		\vspace{0.05cm}
		\stepbox{\scriptsize
			\keypoint{Correct Answer:} Receiving, analyzing, and disseminating suspicious transaction reports (STRs) to law enforcement and other competent authorities.
		}
		\vspace{0.05cm}
		\scriptsize
		\hardconcept{Core Concept:} FIU functions:
		\par \textbf{Receive STRs}: From financial institutions and DNFBPs.
		\par \textbf{Analyze}: Identify patterns, trends, and suspicious activity.
		\par \textbf{Disseminate}: Share intelligence with law enforcement and other FIUs.
		\par \textbf{Feedback}: Provide feedback to reporting entities.
		\vspace{0.03cm}
		\wrongbox{\scriptsize \textbf{Common Trap:} "FIUs investigate and prosecute financial crimes" — INCORRECT. FIUs are \textbf{intelligence units}, not investigative or prosecutorial bodies.}
		\vspace{0.03cm}
		\tipbox{\scriptsize \textbf{Key Insight:} The Egmont Group connects FIUs globally, enabling secure information sharing for cross-border investigations.}
	\end{frame}
	
	\begin{frame}
		\frametitle{\fontsize{11}{13}\selectfont Q13: Law Enforcement — Role}
		\begin{block}{\fontsize{8}{10}\selectfont Topic: Law Enforcement}
			\scriptsize
			What is the role of law enforcement in AML/CFT?
		\end{block}
		\vspace{0.05cm}
		\stepbox{\scriptsize
			\keypoint{Correct Answer:} Investigating, prosecuting, and recovering assets from money laundering and terrorist financing crimes.
		}
		\vspace{0.05cm}
		\scriptsize
		\hardconcept{Core Concept:} Law enforcement:
		\par \textbf{Investigate}: Conduct financial crime investigations.
		\par \textbf{Prosecute}: Bring cases to court.
		\par \textbf{Asset Recovery}: Seize and recover criminal proceeds.
		\par \textbf{International Cooperation}: Work with foreign counterparts.
		\vspace{0.03cm}
		\tipbox{\scriptsize \textbf{Key Insight:} Key agencies include FBI, NCA, Europol, and Interpol. They work closely with FIUs and regulators.}
	\end{frame}
	
	\begin{frame}
		\frametitle{\fontsize{11}{13}\selectfont Q14: Three Pillars Relationship}
		\begin{block}{\fontsize{8}{10}\selectfont Topic: Three Pillars Relationship}
			\scriptsize
			Which statement best describes the relationship between regulators, FIUs, and law enforcement?
		\end{block}
		\vspace{0.05cm}
		\stepbox{\scriptsize
			\keypoint{Correct Answer:} Regulators supervise and enforce compliance, FIUs receive and analyze STRs, and law enforcement investigates and prosecutes — each plays a distinct but complementary role.
		}
		\vspace{0.05cm}
		\scriptsize
		\hardconcept{Core Concept:} The three pillars:
		\par \textbf{Regulators}: Ensure compliance through supervision and enforcement.
		\par \textbf{FIUs}: Produce intelligence from STRs.
		\par \textbf{Law Enforcement}: Use intelligence to investigate and prosecute.
		\vspace{0.03cm}
		\tipbox{\scriptsize \textbf{Key Insight:} Information flows from financial institutions → FIUs → law enforcement. Regulators oversee the institutions.}
	\end{frame}
	
	\begin{frame}
		\frametitle{\fontsize{11}{13}\selectfont Q15: Egmont Group}
		\begin{block}{\fontsize{8}{10}\selectfont Topic: Egmont Group}
			\scriptsize
			What is the primary purpose of the Egmont Group?
		\end{block}
		\vspace{0.05cm}
		\stepbox{\scriptsize
			\keypoint{Correct Answer:} To provide a global network for FIUs to securely share information and cooperate on cross-border financial crime investigations.
		}
		\vspace{0.05cm}
		\scriptsize
		\hardconcept{Core Concept:} Egmont Group:
		\par \textbf{Secure Portal}: Enables confidential information sharing.
		\par \textbf{Typology Reports}: Shared analysis of emerging threats.
		\par \textbf{Capacity Building}: Training and support for member FIUs.
		\par \textbf{Membership}: 160+ FIUs worldwide.
		\vspace{0.03cm}
		\wrongbox{\scriptsize \textbf{Common Trap:} "Egmont Group membership includes automatic FATF evaluation inclusion" — INCORRECT. Egmont membership and FATF evaluations are separate.}
		\vspace{0.03cm}
		\tipbox{\scriptsize \textbf{Key Insight:} The Egmont Group's \textbf{secure network} is essential for cross-border FIU cooperation. It has a Code of Conduct for information sharing.}
	\end{frame}
	
	\begin{frame}
		\frametitle{\fontsize{11}{13}\selectfont Q16: FIU Feedback to Reporting Entities}
		\begin{block}{\fontsize{8}{10}\selectfont Topic: FIU Feedback}
			\scriptsize
			Why is FIU feedback to reporting entities important?
		\end{block}
		\vspace{0.05cm}
		\stepbox{\scriptsize
			\keypoint{Correct Answer:} Feedback helps institutions improve their AML programs by understanding which STRs were valuable and what typologies are emerging.
		}
		\vspace{0.05cm}
		\scriptsize
		\hardconcept{Core Concept:} Feedback benefits:
		\par \textbf{Improves STR Quality}: Institutions learn what is useful.
		\par \textbf{Identifies Typologies}: Highlights emerging ML/TF methods.
		\par \textbf{Builds Trust}: Strengthens the relationship between FIUs and reporting entities.
		\par \textbf{Compliance}: Under \textbf{FATF Recommendation 34}, guidance and feedback are required.
		\vspace{0.03cm}
		\wrongbox{\scriptsize \textbf{Common Trap:} "STR filing is enough; feedback is optional" — INCORRECT. FATF Recommendation 34 requires competent authorities to provide guidance and feedback.}
		\vspace{0.03cm}
		\tipbox{\scriptsize \textbf{Key Insight:} Lack of feedback is a common deficiency in mutual evaluations. It is a violation of FATF Recommendation 34.}
	\end{frame}
	
	% ================================================================
	% SLIDE 19-22: PUBLIC VS PRIVATE SECTOR
	% ================================================================
	\begin{frame}
		\frametitle{\fontsize{11}{13}\selectfont Q17: Public Sector Objectives}
		\begin{block}{\fontsize{8}{10}\selectfont Topic: Public Sector Objectives}
			\scriptsize
			What are the primary objectives of the public sector in AML/CFT?
		\end{block}
		\vspace{0.05cm}
		\stepbox{\scriptsize
			\keypoint{Correct Answer:} Protect national security, maintain financial system integrity, enforce laws, and prosecute criminals.
		}
		\vspace{0.05cm}
		\scriptsize
		\hardconcept{Core Concept:} Public sector objectives:
		\par \textbf{National Security}: Prevent terrorist financing.
		\par \textbf{Financial Integrity}: Protect the financial system from abuse.
		\par \textbf{Law Enforcement}: Investigate and prosecute.
		\par \textbf{Asset Recovery}: Seize criminal proceeds.
		\vspace{0.03cm}
		\tipbox{\scriptsize \textbf{Key Insight:} The public sector's primary motivation is \textbf{public protection}, not commercial interests.}
	\end{frame}
	
	\begin{frame}
		\frametitle{\fontsize{11}{13}\selectfont Q18: Private Sector Objectives}
		\begin{block}{\fontsize{8}{10}\selectfont Topic: Private Sector Objectives}
			\scriptsize
			What are the primary objectives of the private sector in AML/CFT?
		\end{block}
		\vspace{0.05cm}
		\stepbox{\scriptsize
			\keypoint{Correct Answer:} Comply with regulations, protect the institution from financial and reputational damage, manage risk, and detect suspicious activity.
		}
		\vspace{0.05cm}
		\scriptsize
		\hardconcept{Core Concept:} Private sector objectives:
		\par \textbf{Compliance}: Meet regulatory requirements.
		\par \textbf{Risk Management}: Protect against losses and reputational damage.
		\par \textbf{Detection}: Identify and report suspicious activity.
		\par \textbf{Protection}: Safeguard the institution from criminal exposure.
		\vspace{0.03cm}
		\tipbox{\scriptsize \textbf{Key Insight:} The private sector's motivation includes \textbf{commercial risk} and \textbf{regulatory compliance}, not just altruism.}
	\end{frame}
	
	\begin{frame}
		\frametitle{\fontsize{11}{13}\selectfont Q19: Aligning Public and Private Objectives}
		\begin{block}{\fontsize{8}{10}\selectfont Topic: Aligning Public and Private Objectives}
			\scriptsize
			Which of the following best aligns public and private sector objectives in AML/CFT?
		\end{block}
		\vspace{0.05cm}
		\stepbox{\scriptsize
			\keypoint{Correct Answer:} Public-private partnerships (PPPs) that enable information sharing while protecting confidentiality and legal rights.
		}
		\vspace{0.05cm}
		\scriptsize
		\hardconcept{Core Concept:} PPPs:
		\par \textbf{Enable Information Sharing}: Public and private share intelligence.
		\par \textbf{Joint Targeting}: Combine resources against common threats.
		\par \textbf{Protect Rights}: Confidentiality and legal safeguards.
		\vspace{0.03cm}
		\tipbox{\scriptsize \textbf{Key Insight:} PPPs bridge the gap between public security objectives and private risk management.}
	\end{frame}
	
	\begin{frame}
		\frametitle{\fontsize{11}{13}\selectfont Q20: NGO Contributions to AFC}
		\begin{block}{\fontsize{8}{10}\selectfont Topic: NGO Contributions}
			\scriptsize
			Which of the following is a key contribution of NGOs to the global fight against financial crime?
		\end{block}
		\vspace{0.05cm}
		\stepbox{\scriptsize
			\keypoint{Correct Answer:} Publishing research, benchmarks, and indices that inform risk assessments and due diligence practices.
		}
		\vspace{0.05cm}
		\scriptsize
		\hardconcept{Core Concept:} Key NGO contributions:
		\par \textbf{Transparency International}: Corruption Perceptions Index (CPI).
		\par \textbf{Tax Justice Network}: Financial Secrecy Index (FSI).
		\par \textbf{Basel Institute on Governance}: Basel AML Index, ICAR.
		\par \textbf{Global Witness}: Investigative reports on corruption.
		\vspace{0.03cm}
		\wrongbox{\scriptsize \textbf{Common Trap:} "NGOs are the primary enforcers of international AML standards" — INCORRECT. NGOs are \textbf{advocates and researchers}; enforcement is done by governments and regulators.}
		\vspace{0.03cm}
		\tipbox{\scriptsize \textbf{Key Insight:} NGOs provide \textbf{independent, objective data} that complements regulatory information.}
	\end{frame}
	
	% ================================================================
	% SLIDE 23-26: REGULATORY COOPERATION
	% ================================================================
	\begin{frame}
		\frametitle{\fontsize{11}{13}\selectfont Q21: Regulatory Cooperation Mechanisms}
		\begin{block}{\fontsize{8}{10}\selectfont Topic: Regulatory Cooperation}
			\scriptsize
			How do regulatory authorities cooperate on cross-border financial crime investigations?
		\end{block}
		\vspace{0.05cm}
		\stepbox{\scriptsize
			\keypoint{Correct Answer:} Through Mutual Legal Assistance Treaties (MLATs), the Egmont Group network, and coordinated examinations between home and host supervisors.
		}
		\vspace{0.05cm}
		\scriptsize
		\hardconcept{Core Concept:} Key cooperation mechanisms:
		\par \textbf{MLATs}: Formal treaties enabling exchange of evidence.
		\par \textbf{Egmont Group}: Secure information sharing between FIUs.
		\par \textbf{Home-Host Supervision}: Coordinated regulatory examinations.
		\par \textbf{Joint Reviews}: Regulators from different jurisdictions conducting coordinated exams.
		\vspace{0.03cm}
		\tipbox{\scriptsize \textbf{Key Insight:} The \textbf{Basel Committee's Consolidated Supervision} framework requires home supervisors to oversee cross-border institutions with host supervisors coordinating examinations.}
	\end{frame}
	
	\begin{frame}
		\frametitle{\fontsize{11}{13}\selectfont Q22: Mutual Legal Assistance Treaties (MLATs)}
		\begin{block}{\fontsize{8}{10}\selectfont Topic: MLATs}
			\scriptsize
			What is the primary purpose of a Mutual Legal Assistance Treaty (MLAT)?
		\end{block}
		\vspace{0.05cm}
		\stepbox{\scriptsize
			\keypoint{Correct Answer:} To facilitate the exchange of evidence and information between countries for criminal investigations and prosecutions.
		}
		\vspace{0.05cm}
		\scriptsize
		\hardconcept{Core Concept:} MLATs:
		\par \textbf{Evidence Exchange}: Share documents, witness testimony, and records.
		\par \textbf{Asset Freezing}: Freeze assets across borders.
		\par \textbf{Extradition}: Support extradition requests.
		\par \textbf{Legal Framework}: Provide legal basis for cooperation.
		\vspace{0.03cm}
		\wrongbox{\scriptsize \textbf{Common Trap:} "MLATs allow unlimited information sharing without restrictions" — INCORRECT. MLATs are subject to \textbf{legal and procedural requirements}.}
		\vspace{0.03cm}
		\tipbox{\scriptsize \textbf{Key Insight:} \textbf{FATF Recommendation 36} requires countries to provide the widest possible range of mutual legal assistance.}
	\end{frame}
	
	\begin{frame}
		\frametitle{\fontsize{11}{13}\selectfont Q23: Home-Host Supervision}
		\begin{block}{\fontsize{8}{10}\selectfont Topic: Home-Host Supervision}
			\scriptsize
			What is the principle of home-host supervision in AML/CFT?
		\end{block}
		\vspace{0.05cm}
		\stepbox{\scriptsize
			\keypoint{Correct Answer:} The home supervisor has primary responsibility for supervising the institution's group-wide AML program, while host supervisors focus on local compliance.
		}
		\vspace{0.05cm}
		\scriptsize
		\hardconcept{Core Concept:}
		\par \textbf{Home Supervisor}: Primary oversight of group-wide compliance.
		\par \textbf{Host Supervisor}: Local compliance in each jurisdiction.
		\par \textbf{Coordination}: Both supervisors coordinate to avoid duplication.
		\vspace{0.03cm}
		\tipbox{\scriptsize \textbf{Key Insight:} Under \textbf{FATF Recommendation 18}, group-wide AML programs must apply to all branches and subsidiaries. The home supervisor ensures this.}
	\end{frame}
	
	\begin{frame}
		\frametitle{\fontsize{11}{13}\selectfont Q24: Cross-Border Investigation Challenges}
		\begin{block}{\fontsize{8}{10}\selectfont Topic: Cross-Border Challenges}
			\scriptsize
			What is the most significant challenge in cross-border money laundering investigations?
		\end{block}
		\vspace{0.05cm}
		\stepbox{\scriptsize
			\keypoint{Correct Answer:} Different legal frameworks, data protection laws, and procedural requirements across jurisdictions that delay or prevent information sharing.
		}
		\vspace{0.05cm}
		\scriptsize
		\hardconcept{Core Concept:} Challenges:
		\par \textbf{Legal Divergence}: Different ML/TF definitions.
		\par \textbf{Data Protection}: GDPR vs US privacy laws.
		\par \textbf{Procedural Delays}: MLATs can take months or years.
		\par \textbf{Lack of Trust}: Some countries do not cooperate.
		\vspace{0.03cm}
		\tipbox{\scriptsize \textbf{Key Insight:} \textbf{FATF Recommendation 2} requires countries to cooperate internationally. However, legal divergence remains a significant obstacle.}
	\end{frame}
	
	% ================================================================
	% SLIDE 27-30: LEVERAGING REGULATORY REPORTS
	% ================================================================
	\begin{frame}
		\frametitle{\fontsize{11}{13}\selectfont Q25: Using National Risk Assessments}
		\begin{block}{\fontsize{8}{10}\selectfont Topic: Leveraging Regulatory Reports}
			\scriptsize
			How should a compliance team use a National Risk Assessment (NRA)?
		\end{block}
		\vspace{0.05cm}
		\stepbox{\scriptsize
			\keypoint{Correct Answer:} Use the NRA to identify country-level ML/TF threats and vulnerabilities that inform the institution's enterprise-wide risk assessment.
		}
		\vspace{0.03cm}
		\tipbox{\scriptsize \textbf{Key Insight:} Under \textbf{FATF Recommendation 1}, institutions must understand national risks.}
	\end{frame}
	
	\begin{frame}
		\frametitle{\fontsize{11}{13}\selectfont Q26: FIU Typology Reports}
		\begin{block}{\fontsize{8}{10}\selectfont Topic: FIU Typology Reports}
			\scriptsize
			What is the value of FIU typology reports for financial institutions?
		\end{block}
		\vspace{0.05cm}
		\stepbox{\scriptsize
			\keypoint{Correct Answer:} They identify emerging ML/TF methods, enabling institutions to update monitoring scenarios and training programs.
		}
		\vspace{0.03cm}
		\tipbox{\scriptsize \textbf{Key Insight:} Typology reports are \textbf{actionable intelligence} that help institutions stay ahead of criminal innovation.}
	\end{frame}
	
	\begin{frame}
		\frametitle{\fontsize{11}{13}\selectfont Q27: Regulatory Enforcement Actions}
		\begin{block}{\fontsize{8}{10}\selectfont Topic: Regulatory Enforcement Actions}
			\scriptsize
			How can a compliance team learn from regulatory enforcement actions?
		\end{block}
		\vspace{0.05cm}
		\stepbox{\scriptsize
			\keypoint{Correct Answer:} Review the deficiencies identified in enforcement actions to identify gaps in their own program and implement corrective measures.
		}
		\vspace{0.03cm}
		\tipbox{\scriptsize \textbf{Key Insight:} Enforcement actions provide a \textbf{roadmap of what regulators consider deficient}.}
	\end{frame}
	
	\begin{frame}
		\frametitle{\fontsize{11}{13}\selectfont Q28: Sectoral Risk Assessments}
		\begin{block}{\fontsize{8}{10}\selectfont Topic: Sectoral Risk Assessments}
			\scriptsize
			What is the purpose of sectoral risk assessments?
		\end{block}
		\vspace{0.05cm}
		\stepbox{\scriptsize
			\keypoint{Correct Answer:} To identify ML/TF risks specific to industry sectors (e.g., banking, real estate, gambling) and inform sector-specific controls.
		}
		\vspace{0.03cm}
		\tipbox{\scriptsize \textbf{Key Insight:} Sectoral risk assessments complement NRAs by providing \textbf{industry-specific} risk information.}
	\end{frame}
	
	% ================================================================
	% SLIDE 31-34: NON-GOVERNMENT GUIDANCE
	% ================================================================
	\begin{frame}
		\frametitle{\fontsize{11}{13}\selectfont Q29: Wolfsberg Group Guidance}
		\begin{block}{\fontsize{8}{10}\selectfont Topic: Wolfsberg Group}
			\scriptsize
			What is the primary purpose of the Wolfsberg Group's AML guidance?
		\end{block}
		\vspace{0.05cm}
		\stepbox{\scriptsize
			\keypoint{Correct Answer:} Providing industry best practices for AML/CFT, including the Correspondent Banking Due Diligence Questionnaire (CBDDQ).
		}
		\vspace{0.03cm}
		\tipbox{\scriptsize \textbf{Key Insight:} The Wolfsberg Group consists of 13 global banks. Its guidance is widely adopted as \textbf{industry best practice}.}
	\end{frame}
	
	\begin{frame}
		\frametitle{\fontsize{11}{13}\selectfont Q30: Basel Committee Guidance}
		\begin{block}{\fontsize{8}{10}\selectfont Topic: Basel Committee}
			\scriptsize
			What does the Basel Committee on Banking Supervision provide for AML/CFT?
		\end{block}
		\vspace{0.05cm}
		\stepbox{\scriptsize
			\keypoint{Correct Answer:} Principles for sound management of AML/CFT risks, including the Three Lines of Defense model and consolidated supervision.
		}
		\vspace{0.03cm}
		\tipbox{\scriptsize \textbf{Key Insight:} The Basel Committee's guidance is referenced by regulators globally and forms the basis of many national frameworks.}
	\end{frame}
	
	\begin{frame}
		\frametitle{\fontsize{11}{13}\selectfont Q31: ICAR Asset Recovery Support}
		\begin{block}{\fontsize{8}{10}\selectfont Topic: ICAR}
			\scriptsize
			What does the Basel Institute's International Centre for Asset Recovery (ICAR) provide?
		\end{block}
		\vspace{0.05cm}
		\stepbox{\scriptsize
			\keypoint{Correct Answer:} Case mentoring and facilitating international cooperation on asset recovery.
		}
		\vspace{0.03cm}
		\tipbox{\scriptsize \textbf{Key Insight:} ICAR has supported cases in the Balkans, Latin America, and Asia-Pacific, recovering billions in stolen assets.}
	\end{frame}
	
	\begin{frame}
		\frametitle{\fontsize{11}{13}\selectfont Q32: Using CPI for Due Diligence}
		\begin{block}{\fontsize{8}{10}\selectfont Topic: Corruption Perceptions Index}
			\scriptsize
			How should a compliance team use Transparency International's Corruption Perceptions Index (CPI)?
		\end{block}
		\vspace{0.05cm}
		\stepbox{\scriptsize
			\keypoint{Correct Answer:} Use CPI scores to assess corruption risk in jurisdictions and apply enhanced due diligence to clients from low-scoring countries.
		}
		\vspace{0.03cm}
		\tipbox{\scriptsize \textbf{Key Insight:} A CPI score below 50 should trigger enhanced scrutiny in country risk assessments.}
	\end{frame}
	
	% ================================================================
	% SLIDE 35-38: NATIONAL AND SECTORAL RISK ASSESSMENTS
	% ================================================================
	\begin{frame}
		\frametitle{\fontsize{11}{13}\selectfont Q33: Purpose of NRAs}
		\begin{block}{\fontsize{8}{10}\selectfont Topic: National Risk Assessments}
			\scriptsize
			What is the primary purpose of a National Risk Assessment (NRA)?
		\end{block}
		\vspace{0.05cm}
		\stepbox{\scriptsize
			\keypoint{Correct Answer:} To identify, assess, and understand the ML/TF risks facing a country to inform policy and resource allocation.
		}
		\vspace{0.03cm}
		\tipbox{\scriptsize \textbf{Key Insight:} Under \textbf{FATF Recommendation 1}, countries must conduct NRAs. They are the foundation of the risk-based approach.}
	\end{frame}
	
	\begin{frame}
		\frametitle{\fontsize{11}{13}\selectfont Q34: NRA Components}
		\begin{block}{\fontsize{8}{10}\selectfont Topic: NRA Components}
			\scriptsize
			What are the three components of an NRA?
		\end{block}
		\vspace{0.05cm}
		\stepbox{\scriptsize
			\keypoint{Correct Answer:} Threat assessment, vulnerability assessment, and consequence assessment.
		}
		\vspace{0.03cm}
		\tipbox{\scriptsize \textbf{Key Insight:} \textbf{Threat} = who is committing crime; \textbf{Vulnerability} = where systems are weak; \textbf{Consequence} = impact of crime.}
	\end{frame}
	
	\begin{frame}
		\frametitle{\fontsize{11}{13}\selectfont Q35: Sectoral Risk Assessments}
		\begin{block}{\fontsize{8}{10}\selectfont Topic: Sectoral Risk Assessments}
			\scriptsize
			What is the purpose of sectoral risk assessments?
		\end{block}
		\vspace{0.05cm}
		\stepbox{\scriptsize
			\keypoint{Correct Answer:} To identify ML/TF risks specific to industry sectors and inform sector-specific controls.
		}
		\vspace{0.03cm}
		\tipbox{\scriptsize \textbf{Key Insight:} Sectoral risk assessments are used by regulators to prioritize supervisory resources.}
	\end{frame}
	
	\begin{frame}
		\frametitle{\fontsize{11}{13}\selectfont Q36: Using NRAs in EWRA}
		\begin{block}{\fontsize{8}{10}\selectfont Topic: NRA and EWRA}
			\scriptsize
			How does an NRA impact an institution's enterprise-wide risk assessment (EWRA)?
		\end{block}
		\vspace{0.05cm}
		\stepbox{\scriptsize
			\keypoint{Correct Answer:} The NRA provides the country-level risk context that must be incorporated into the EWRA, ensuring the institution's risk assessment reflects national threats.
		}
		\vspace{0.03cm}
		\tipbox{\scriptsize \textbf{Key Insight:} Under \textbf{FATF Recommendation 1}, institutions must understand national risks and incorporate them into their EWRA.}
	\end{frame}
	
	% ================================================================
	% SLIDE 39-44: US AND EU REGULATIONS
	% ================================================================
	\begin{frame}
		\frametitle{\fontsize{11}{13}\selectfont Q37: US BSA Framework}
		\begin{block}{\fontsize{8}{10}\selectfont Topic: US BSA}
			\scriptsize
			What is the foundation of US AML regulation?
		\end{block}
		\vspace{0.05cm}
		\stepbox{\scriptsize
			\keypoint{Correct Answer:} The Bank Secrecy Act (BSA) of 1970, as amended by the PATRIOT Act of 2001 and subsequent legislation.
		}
		\vspace{0.03cm}
		\tipbox{\scriptsize \textbf{Key Insight:} The BSA establishes the \textbf{five pillars}: internal controls, compliance officer, training, independent testing, and CDD (including BO).}
	\end{frame}
	
	\begin{frame}
		\frametitle{\fontsize{11}{13}\selectfont Q38: US PATRIOT Act Section 311}
		\begin{block}{\fontsize{8}{10}\selectfont Topic: PATRIOT Act Section 311}
			\scriptsize
			What does Section 311 of the PATRIOT Act allow?
		\end{block}
		\vspace{0.05cm}
		\stepbox{\scriptsize
			\keypoint{Correct Answer:} The Treasury Department to designate foreign jurisdictions, institutions, or transactions as "primary money laundering concerns" and impose special measures.
		}
		\vspace{0.03cm}
		\tipbox{\scriptsize \textbf{Key Insight:} Section 311 is a powerful tool that can require enhanced due diligence or prohibit relationships entirely.}
	\end{frame}
	
	\begin{frame}
		\frametitle{\fontsize{11}{13}\selectfont Q39: US OFAC Sanctions}
		\begin{block}{\fontsize{8}{10}\selectfont Topic: OFAC}
			\scriptsize
			What is the role of the US Office of Foreign Assets Control (OFAC)?
		\end{block}
		\vspace{0.05cm}
		\stepbox{\scriptsize
			\keypoint{Correct Answer:} Administering and enforcing US economic sanctions against designated countries, individuals, and entities.
		}
		\vspace{0.03cm}
		\tipbox{\scriptsize \textbf{Key Insight:} OFAC sanctions are \textbf{strict liability} — no intent is required for a violation.}
	\end{frame}
	
	\begin{frame}
		\frametitle{\fontsize{11}{13}\selectfont Q40: EU AML Directives}
		\begin{block}{\fontsize{8}{10}\selectfont Topic: EU AML Directives}
			\scriptsize
			What is the current framework for EU AML regulation?
		\end{block}
		\vspace{0.05cm}
		\stepbox{\scriptsize
			\keypoint{Correct Answer:} The AML Directives (4AMLD, 5AMLD, 6AMLD) transposed into national law, with the AMLR and AMLA providing centralized oversight.
		}
		\vspace{0.03cm}
		\tipbox{\scriptsize \textbf{Key Insight:} 6AMLD expanded predicate offenses and criminalized aiding and abetting. AMLA (Frankfurt, 2026) provides direct supervision.}
	\end{frame}
	
	\begin{frame}
		\frametitle{\fontsize{11}{13}\selectfont Q41: 6AMLD Key Provisions}
		\begin{block}{\fontsize{8}{10}\selectfont Topic: 6AMLD}
			\scriptsize
			What are key provisions of the 6th Anti-Money Laundering Directive (6AMLD)?
		\end{block}
		\vspace{0.05cm}
		\stepbox{\scriptsize
			\keypoint{Correct Answer:} Expanded predicate offenses, criminalization of aiding and abetting, minimum 4-year imprisonment for ML offenses, and inclusion of soccer agents and CASPs.
		}
		\vspace{0.03cm}
		\tipbox{\scriptsize \textbf{Key Insight:} 6AMLD added \textbf{soccer agents} and \textbf{cryptoasset service providers} to the AML framework.}
	\end{frame}
	
	\begin{frame}
		\frametitle{\fontsize{11}{13}\selectfont Q42: AMLA Role}
		\begin{block}{\fontsize{8}{10}\selectfont Topic: AMLA}
			\scriptsize
			What is the role of the Anti-Money Laundering Authority (AMLA)?
		\end{block}
		\vspace{0.05cm}
		\stepbox{\scriptsize
			\keypoint{Correct Answer:} Centralized EU supervisor with direct powers over highest-risk cross-border entities, including on-site inspections, fines, and license revocation.
		}
		\vspace{0.03cm}
		\tipbox{\scriptsize \textbf{Key Insight:} AMLA (Frankfurt, 2026) provides direct oversight of entities with cross-border operations and significant risk profiles.}
	\end{frame}
	
	% ================================================================
	% SLIDE 45-48: JURISDICTIONAL AWARENESS
	% ================================================================
	\begin{frame}
		\frametitle{\fontsize{11}{13}\selectfont Q43: Why Jurisdictional Awareness Matters}
		\begin{block}{\fontsize{8}{10}\selectfont Topic: Jurisdictional Awareness}
			\scriptsize
			Why is it important to be aware of AFC regulations in different jurisdictions?
		\end{block}
		\vspace{0.05cm}
		\stepbox{\scriptsize
			\keypoint{Correct Answer:} Different jurisdictions have different AML/CFT requirements, sanctions regimes, and enforcement approaches, creating compliance gaps and conflicts.
		}
		\vspace{0.03cm}
		\tipbox{\scriptsize \textbf{Key Insight:} \textbf{FATF Recommendation 18} requires group-wide AML programs that apply to all branches and subsidiaries.}
	\end{frame}
	
	\begin{frame}
		\frametitle{\fontsize{11}{13}\selectfont Q44: Regulatory Divergence Example}
		\begin{block}{\fontsize{8}{10}\selectfont Topic: Regulatory Divergence}
			\scriptsize
			Which of the following is an example of regulatory divergence between jurisdictions?
		\end{block}
		\vspace{0.05cm}
		\stepbox{\scriptsize
			\keypoint{Correct Answer:} The UK Bribery Act prohibits facilitation payments, while the US FCPA allows them in certain circumstances.
		}
		\vspace{0.03cm}
		\tipbox{\scriptsize \textbf{Key Insight:} This divergence creates compliance challenges for global institutions operating in both jurisdictions.}
	\end{frame}
	
	\begin{frame}
		\frametitle{\fontsize{11}{13}\selectfont Q45: Managing Jurisdictional Conflicts}
		\begin{block}{\fontsize{8}{10}\selectfont Topic: Managing Conflicts}
			\scriptsize
			How should a global institution manage conflicts between jurisdictional AML requirements?
		\end{block}
		\vspace{0.05cm}
		\stepbox{\scriptsize
			\keypoint{Correct Answer:} Apply the stricter standard where there is a conflict, and seek legal advice on data protection and information-sharing restrictions.
		}
		\vspace{0.03cm}
		\tipbox{\scriptsize \textbf{Key Insight:} The \textbf{principle of strictest standard} is commonly applied when regulatory requirements conflict.}
	\end{frame}
	
	\begin{frame}
		\frametitle{\fontsize{11}{13}\selectfont Q46: Sanctions Differences}
		\begin{block}{\fontsize{8}{10}\selectfont Topic: Sanctions Differences}
			\scriptsize
			What is a common challenge with sanctions compliance across jurisdictions?
		\end{block}
		\vspace{0.05cm}
		\stepbox{\scriptsize
			\keypoint{Correct Answer:} US OFAC sanctions may differ from EU or UN sanctions, requiring institutions to comply with multiple lists simultaneously.
		}
		\vspace{0.03cm}
		\tipbox{\scriptsize \textbf{Key Insight:} Global institutions must screen against \textbf{all applicable sanctions lists} (UN, US, EU, and local).}
	\end{frame}
	
	% ================================================================
	% SLIDE 49-54: PUBLIC-PRIVATE PARTNERSHIPS
	% ================================================================
	\begin{frame}
		\frametitle{\fontsize{11}{13}\selectfont Q47: PPP Importance}
		\begin{block}{\fontsize{8}{10}\selectfont Topic: PPP Importance}
			\scriptsize
			Why are public-private partnerships important in AML/CFT?
		\end{block}
		\vspace{0.05cm}
		\stepbox{\scriptsize
			\keypoint{Correct Answer:} They enable timely information sharing between public and private sectors, allowing joint targeting of high-risk actors.
		}
		\vspace{0.03cm}
		\tipbox{\scriptsize \textbf{Key Insight:} Successful PPP models include UK JMLIT, Australia Fintel, and Singapore ACIP.}
	\end{frame}
	
	\begin{frame}
		\frametitle{\fontsize{11}{13}\selectfont Q48: PPP Goals}
		\begin{block}{\fontsize{8}{10}\selectfont Topic: PPP Goals}
			\scriptsize
			What do PPPs aim to achieve in AML/CFT?
		\end{block}
		\vspace{0.05cm}
		\stepbox{\scriptsize
			\keypoint{Correct Answer:} Real-time intelligence sharing, joint targeting of threats, and disruption of criminal networks.
		}
		\vspace{0.03cm}
		\tipbox{\scriptsize \textbf{Key Insight:} PPPs aim to \textbf{disrupt} financial crime, not just detect it after the fact.}
	\end{frame}
	
	\begin{frame}
		\frametitle{\fontsize{11}{13}\selectfont Q49: PPP Information Sharing}
		\begin{block}{\fontsize{8}{10}\selectfont Topic: PPP Information Sharing}
			\scriptsize
			What type of information is typically shared in a PPP?
		\end{block}
		\vspace{0.05cm}
		\stepbox{\scriptsize
			\keypoint{Correct Answer:} Threat intelligence, typology information, and specific indicators of suspicious activity, subject to confidentiality and data protection safeguards.
		}
		\vspace{0.03cm}
		\tipbox{\scriptsize \textbf{Key Insight:} Information sharing in PPPs is \textbf{protected} and \textbf{confidential}.}
	\end{frame}
	
	\begin{frame}
		\frametitle{\fontsize{11}{13}\selectfont Q50: PPP Real-Time Intelligence}
		\begin{block}{\fontsize{8}{10}\selectfont Topic: PPP Real-Time Intelligence}
			\scriptsize
			Which approach best aligns with the goals of a PPP framework?
		\end{block}
		\vspace{0.05cm}
		\stepbox{\scriptsize
			\keypoint{Correct Answer:} Establishing a formal platform where regulators and institutions engage in direct two-way communication on specific threats.
		}
		\vspace{0.03cm}
		\tipbox{\scriptsize \textbf{Key Insight:} Anonymized, after-the-fact reporting is NOT sufficient for disruption. Real-time, two-way communication is essential.}
	\end{frame}
	
	\begin{frame}
		\frametitle{\fontsize{11}{13}\selectfont Q51: PPP Strengths}
		\begin{block}{\fontsize{8}{10}\selectfont Topic: PPP Strengths}
			\scriptsize
			What is a strength of public-private partnerships?
		\end{block}
		\vspace{0.05cm}
		\stepbox{\scriptsize
			\keypoint{Correct Answer:} Combining public sector intelligence with private sector data to create a comprehensive threat picture.
		}
		\vspace{0.03cm}
		\tipbox{\scriptsize \textbf{Key Insight:} PPPs leverage \textbf{complementary capabilities} — public has legal powers, private has data.}
	\end{frame}
	
	\begin{frame}
		\frametitle{\fontsize{11}{13}\selectfont Q52: PPP Limitations}
		\begin{block}{\fontsize{8}{10}\selectfont Topic: PPP Limitations}
			\scriptsize
			What is a limitation of public-private partnerships?
		\end{block}
		\vspace{0.05cm}
		\stepbox{\scriptsize
			\keypoint{Correct Answer:} Data protection laws and confidentiality obligations may restrict the type and amount of information that can be shared.
		}
		\vspace{0.03cm}
		\tipbox{\scriptsize \textbf{Key Insight:} GDPR, legal privilege, and tipping-off restrictions limit PPP information sharing.}
	\end{frame}
	
	\begin{frame}
		\frametitle{\fontsize{11}{13}\selectfont Q53: PPP Resource Requirements}
		\begin{block}{\fontsize{8}{10}\selectfont Topic: PPP Resources}
			\scriptsize
			What is a resource requirement for effective PPPs?
		\end{block}
		\vspace{0.05cm}
		\stepbox{\scriptsize
			\keypoint{Correct Answer:} Dedicated personnel, technology infrastructure, and ongoing commitment from both public and private sectors.
		}
		\vspace{0.03cm}
		\tipbox{\scriptsize \textbf{Key Insight:} PPPs require \textbf{sustained investment} in people, systems, and trust-building.}
	\end{frame}
	
	\begin{frame}
		\frametitle{\fontsize{11}{13}\selectfont Q54: PPP Trust Building}
		\begin{block}{\fontsize{8}{10}\selectfont Topic: PPP Trust}
			\scriptsize
			Why is trust essential for effective PPPs?
		\end{block}
		\vspace{0.05cm}
		\stepbox{\scriptsize
			\keypoint{Correct Answer:} Trust enables open information sharing and collaboration, which is essential for identifying and disrupting criminal activity.
		}
		\vspace{0.03cm}
		\tipbox{\scriptsize \textbf{Key Insight:} Trust takes time to build and can be easily broken. Clear protocols and confidentiality safeguards are essential.}
	\end{frame}
	
	% ================================================================
	% SLIDE 55-58: PRIVATE SECTOR COLLABORATION
	% ================================================================
	\begin{frame}
		\frametitle{\fontsize{11}{13}\selectfont Q55: Private Sector Collaboration}
		\begin{block}{\fontsize{8}{10}\selectfont Topic: Private Sector Collaboration}
			\scriptsize
			How can the private sector foster collaboration in fighting financial crime?
		\end{block}
		\vspace{0.05cm}
		\stepbox{\scriptsize
			\keypoint{Correct Answer:} By participating in information-sharing initiatives like PPPs, industry forums, and cross-institution intelligence sharing.
		}
		\vspace{0.03cm}
		\tipbox{\scriptsize \textbf{Key Insight:} Industry forums, such as the Wolfsberg Group and IIF, enable private sector collaboration.}
	\end{frame}
	
	\begin{frame}
		\frametitle{\fontsize{11}{13}\selectfont Q56: Section 314(b) Sharing}
		\begin{block}{\fontsize{8}{10}\selectfont Topic: Section 314(b)}
			\scriptsize
			What does Section 314(b) of the US PATRIOT Act allow?
		\end{block}
		\vspace{0.05cm}
		\stepbox{\scriptsize
			\keypoint{Correct Answer:} Financial institutions to share information with each other for AML purposes if they file an opt-in notice with FinCEN.
		}
		\vspace{0.03cm}
		\tipbox{\scriptsize \textbf{Key Insight:} Sharing without opt-in is a violation. Confidentiality and safe harbor protections apply.}
	\end{frame}
	
	\begin{frame}
		\frametitle{\fontsize{11}{13}\selectfont Q57: Industry Forum Benefits}
		\begin{block}{\fontsize{8}{10}\selectfont Topic: Industry Forums}
			\scriptsize
			What are the benefits of industry forums for AML collaboration?
		\end{block}
		\vspace{0.05cm}
		\stepbox{\scriptsize
			\keypoint{Correct Answer:} Sharing typologies, best practices, and emerging threats across institutions in a sector.
		}
		\vspace{0.03cm}
		\tipbox{\scriptsize \textbf{Key Insight:} Industry forums help institutions stay ahead of emerging risks.}
	\end{frame}
	
	\begin{frame}
		\frametitle{\fontsize{11}{13}\selectfont Q58: Cross-Institution Intelligence Sharing}
		\begin{block}{\fontsize{8}{10}\selectfont Topic: Cross-Institution Sharing}
			\scriptsize
			What should institutions consider when sharing intelligence with other institutions?
		\end{block}
		\vspace{0.05cm}
		\stepbox{\scriptsize
			\keypoint{Correct Answer:} Legal authority, confidentiality, data protection, and the risk of tipping-off subjects.
		}
		\vspace{0.03cm}
		\tipbox{\scriptsize \textbf{Key Insight:} Sharing must be \textbf{legal, proportionate, and protected}.}
	\end{frame}
	
	% ================================================================
	% SLIDE 59-62: DATA SHARING CONSIDERATIONS
	% ================================================================
	\begin{frame}
		\frametitle{\fontsize{11}{13}\selectfont Q59: Data Sharing Legal Authority}
		\begin{block}{\fontsize{8}{10}\selectfont Topic: Data Sharing}
			\scriptsize
			What is the most important consideration before sharing data for AML purposes?
		\end{block}
		\vspace{0.05cm}
		\stepbox{\scriptsize
			\keypoint{Correct Answer:} Legal authority to share the data under applicable laws and regulations.
		}
		\vspace{0.03cm}
		\tipbox{\scriptsize \textbf{Key Insight:} Sharing without legal authority is a violation and can lead to enforcement action.}
	\end{frame}
	
	\begin{frame}
		\frametitle{\fontsize{11}{13}\selectfont Q60: Data Protection and Sharing}
		\begin{block}{\fontsize{8}{10}\selectfont Topic: Data Protection}
			\scriptsize
			How do data protection laws impact AML information sharing?
		\end{block}
		\vspace{0.05cm}
		\stepbox{\scriptsize
			\keypoint{Correct Answer:} They restrict the type and amount of personal data that can be shared, requiring institutions to comply with GDPR or other privacy laws.
		}
		\vspace{0.03cm}
		\tipbox{\scriptsize \textbf{Key Insight:} GDPR Article 49(1)(d) allows transfers for "important public interest" including AML.}
	\end{frame}
	
	\begin{frame}
		\frametitle{\fontsize{11}{13}\selectfont Q61: Tipping-Off Risk}
		\begin{block}{\fontsize{8}{10}\selectfont Topic: Tipping-Off}
			\scriptsize
			What is the risk of tipping-off in AML information sharing?
		\end{block}
		\vspace{0.05cm}
		\stepbox{\scriptsize
			\keypoint{Correct Answer:} Inadvertently alerting the subject of an investigation that they are under scrutiny, allowing them to destroy evidence or flee.
		}
		\vspace{0.03cm}
		\tipbox{\scriptsize \textbf{Key Insight:} Tipping-off is a criminal offense in many jurisdictions. Institutions must have safeguards.}
	\end{frame}
	
	\begin{frame}
		\frametitle{\fontsize{11}{13}\selectfont Q62: Data Quality in Sharing}
		\begin{block}{\fontsize{8}{10}\selectfont Topic: Data Quality}
			\scriptsize
			Why is data quality important when sharing information for AML purposes?
		\end{block}
		\vspace{0.05cm}
		\stepbox{\scriptsize
			\keypoint{Correct Answer:} Inaccurate or incomplete information can lead to false positives, wasted resources, and missed detections.
		}
		\vspace{0.03cm}
		\tipbox{\scriptsize \textbf{Key Insight:} GIGO — Garbage In, Garbage Out — applies to information sharing too.}
	\end{frame}
	
	% ================================================================
	% SLIDE 63-66: DATA PRIVACY
	% ================================================================
	\begin{frame}
		\frametitle{\fontsize{11}{13}\selectfont Q63: GDPR and AML}
		\begin{block}{\fontsize{8}{10}\selectfont Topic: GDPR and AML}
			\scriptsize
			How does GDPR impact AML/CFT applications?
		\end{block}
		\vspace{0.05cm}
		\stepbox{\scriptsize
			\keypoint{Correct Answer:} GDPR limits the collection, use, and sharing of personal data for AML purposes, requiring institutions to balance AML obligations with privacy rights.
		}
		\vspace{0.03cm}
		\tipbox{\scriptsize \textbf{Key Insight:} GDPR Articles 5, 6, and 9 set the framework for AML data processing.}
	\end{frame}
	
	\begin{frame}
		\frametitle{\fontsize{11}{13}\selectfont Q64: Data Minimization}
		\begin{block}{\fontsize{8}{10}\selectfont Topic: Data Minimization}
			\scriptsize
			What is the principle of data minimization in AML?
		\end{block}
		\vspace{0.05cm}
		\stepbox{\scriptsize
			\keypoint{Correct Answer:} Collect only personal data that is adequate, relevant, and limited to what is necessary for AML purposes.
		}
		\vspace{0.03cm}
		\tipbox{\scriptsize \textbf{Key Insight:} Data minimization is a core GDPR principle (Article 5(1)(c)).}
	\end{frame}
	
	\begin{frame}
		\frametitle{\fontsize{11}{13}\selectfont Q65: Purpose Limitation}
		\begin{block}{\fontsize{8}{10}\selectfont Topic: Purpose Limitation}
			\scriptsize
			What is the principle of purpose limitation in AML data processing?
		\end{block}
		\vspace{0.05cm}
		\stepbox{\scriptsize
			\keypoint{Correct Answer:} Personal data must be collected for specified, explicit, and legitimate purposes and not further processed in a way incompatible with those purposes.
		}
		\vspace{0.03cm}
		\tipbox{\scriptsize \textbf{Key Insight:} Purpose limitation is a core GDPR principle (Article 5(1)(b)).}
	\end{frame}
	
	\begin{frame}
		\frametitle{\fontsize{11}{13}\selectfont Q66: Cross-Border Data Transfer}
		\begin{block}{\fontsize{8}{10}\selectfont Topic: Cross-Border Transfer}
			\scriptsize
			What is a challenge for cross-border AML information sharing under GDPR?
		\end{block}
		\vspace{0.03cm}
		\stepbox{\scriptsize
			\keypoint{Correct Answer:} GDPR restricts transfers of personal data to countries without adequate data protection, requiring appropriate safeguards.
		}
		\vspace{0.03cm}
		\tipbox{\scriptsize \textbf{Key Insight:} GDPR Article 49(1)(d) allows transfers for "important public interest" including AML and law enforcement.}
	\end{frame}
	
	% ================================================================
	% SLIDE 67-70: CONSUMER PROTECTION
	% ================================================================
	\begin{frame}
		\frametitle{\fontsize{11}{13}\selectfont Q67: Consumer Protection vs AML}
		\begin{block}{\fontsize{8}{10}\selectfont Topic: Consumer Protection}
			\scriptsize
			What is a tension between consumer protection and AML requirements?
		\end{block}
		\vspace{0.05cm}
		\stepbox{\scriptsize
			\keypoint{Correct Answer:} Consumer protection requires transparency and disclosure, while AML may require secrecy to avoid tipping-off.
		}
		\vspace{0.03cm}
		\tipbox{\scriptsize \textbf{Key Insight:} AML obligations often prevail, but institutions must have policies to handle conflicts.}
	\end{frame}
	
	\begin{frame}
		\frametitle{\fontsize{11}{13}\selectfont Q68: Account Freezing Conflicts}
		\begin{block}{\fontsize{8}{10}\selectfont Topic: Account Freezing}
			\scriptsize
			What is a conflict between consumer protection and AML in account freezing?
		\end{block}
		\vspace{0.05cm}
		\stepbox{\scriptsize
			\keypoint{Correct Answer:} Consumer protection may require notice, while AML may require immediate freezing without notice.
		}
		\vspace{0.03cm}
		\tipbox{\scriptsize \textbf{Key Insight:} AML freezing requirements under FATF and sanctions regimes take precedence.}
	\end{frame}
	
	\begin{frame}
		\frametitle{\fontsize{11}{13}\selectfont Q69: Fair Treatment vs Risk-Based}
		\begin{block}{\fontsize{8}{10}\selectfont Topic: Fair Treatment}
			\scriptsize
			How can a risk-based approach conflict with fair treatment of customers?
		\end{block}
		\vspace{0.05cm}
		\stepbox{\scriptsize
			\keypoint{Correct Answer:} Risk-based AML requires treating high-risk customers differently (EDD), which may be perceived as unfair or discriminatory.
		}
		\vspace{0.03cm}
		\tipbox{\scriptsize \textbf{Key Insight:} Institutions must balance fair treatment with risk-based compliance.}
	\end{frame}
	
	\begin{frame}
		\frametitle{\fontsize{11}{13}\selectfont Q70: Balancing Consumer Protection}
		\begin{block}{\fontsize{8}{10}\selectfont Topic: Balancing}
			\scriptsize
			How should institutions balance consumer protection and AML obligations?
		\end{block}
		\vspace{0.05cm}
		\stepbox{\scriptsize
			\keypoint{Correct Answer:} Develop clear policies that balance both, with legal guidance on handling conflicts, and communicate transparently with customers where allowed.
		}
		\vspace{0.03cm}
		\tipbox{\scriptsize \textbf{Key Insight:} A clear policy framework helps staff navigate conflicts.}
	\end{frame}
	
	% ================================================================
	% SLIDE 71-74: ESG
	% ================================================================
	\begin{frame}
		\frametitle{\fontsize{11}{13}\selectfont Q71: ESG and AML Intersection}
		\begin{block}{\fontsize{8}{10}\selectfont Topic: ESG and AML}
			\scriptsize
			How do ESG factors intersect with AML/CFT compliance?
		\end{block}
		\vspace{0.05cm}
		\stepbox{\scriptsize
			\keypoint{Correct Answer:} ESG factors, particularly governance and social responsibility, are increasingly expected to be considered in AML risk assessments.
		}
		\vspace{0.03cm}
		\tipbox{\scriptsize \textbf{Key Insight:} Regulators expect institutions to consider ESG in AML programs.}
	\end{frame}
	
	\begin{frame}
		\frametitle{\fontsize{11}{13}\selectfont Q72: Governance and AML}
		\begin{block}{\fontsize{8}{10}\selectfont Topic: Governance}
			\scriptsize
			What does the "G" in ESG mean for AML compliance?
		\end{block}
		\vspace{0.05cm}
		\stepbox{\scriptsize
			\keypoint{Correct Answer:} Board oversight, compliance culture, ethics, and transparency in governance structures.
		}
		\vspace{0.03cm}
		\tipbox{\scriptsize \textbf{Key Insight:} Governance is the most directly relevant ESG factor for AML.}
	\end{frame}
	
	\begin{frame}
		\frametitle{\fontsize{11}{13}\selectfont Q73: Environmental Crime and AML}
		\begin{block}{\fontsize{8}{10}\selectfont Topic: Environmental Crime}
			\scriptsize
			Why is environmental crime increasingly linked to money laundering?
		\end{block}
		\vspace{0.05cm}
		\stepbox{\scriptsize
			\keypoint{Correct Answer:} Environmental crimes (illegal logging, mining, wildlife trafficking) generate high-volume cash proceeds that require money laundering.
		}
		\vspace{0.03cm}
		\tipbox{\scriptsize \textbf{Key Insight:} FATF has issued guidance on Environmental Crime and Money Laundering.}
	\end{frame}
	
	\begin{frame}
		\frametitle{\fontsize{11}{13}\selectfont Q74: Social Responsibility and AML}
		\begin{block}{\fontsize{8}{10}\selectfont Topic: Social Responsibility}
			\scriptsize
			What social factors are relevant to AML risk assessment?
		\end{block}
		\vspace{0.05cm}
		\stepbox{\scriptsize
			\keypoint{Correct Answer:} Human trafficking, modern slavery, and exploitation — all require money laundering for proceeds.
		}
		\vspace{0.03cm}
		\tipbox{\scriptsize \textbf{Key Insight:} Social factors like human trafficking are increasingly integrated into AML risk assessments.}
	\end{frame}
	
	% ================================================================
	% SLIDE 75-78: FINANCIAL INCLUSION
	% ================================================================
	\begin{frame}
		\frametitle{\fontsize{11}{13}\selectfont Q75: Financial Inclusion vs AML}
		\begin{block}{\fontsize{8}{10}\selectfont Topic: Financial Inclusion}
			\scriptsize
			What is the primary tension between financial inclusion and AML?
		\end{block}
		\vspace{0.05cm}
		\stepbox{\scriptsize
			\keypoint{Correct Answer:} AML due diligence requirements may create barriers for underserved populations lacking traditional identification.
		}
		\vspace{0.03cm}
		\tipbox{\scriptsize \textbf{Key Insight:} FATF has issued guidance on Financial Inclusion and proportionality.}
	\end{frame}
	
	\begin{frame}
		\frametitle{\fontsize{11}{13}\selectfont Q76: Proportionality in Financial Inclusion}
		\begin{block}{\fontsize{8}{10}\selectfont Topic: Proportionality}
			\scriptsize
			How can institutions balance financial inclusion with AML controls?
		\end{block}
		\vspace{0.05cm}
		\stepbox{\scriptsize
			\keypoint{Correct Answer:} Apply proportionate CDD measures based on risk, using simplified due diligence for lower-risk customers.
		}
		\vspace{0.03cm}
		\tipbox{\scriptsize \textbf{Key Insight:} Simplified due diligence is permitted under \textbf{FATF Recommendation 10} for lower-risk customers.}
	\end{frame}
	
	\begin{frame}
		\frametitle{\fontsize{11}{13}\selectfont Q77: De-risking and Inclusion}
		\begin{block}{\fontsize{8}{10}\selectfont Topic: De-risking}
			\scriptsize
			What is de-risking and how does it impact financial inclusion?
		\end{block}
		\vspace{0.05cm}
		\stepbox{\scriptsize
			\keypoint{Correct Answer:} De-risking is exiting entire market segments to avoid AML risk, which can exclude legitimate customers and undermine financial inclusion.
		}
		\vspace{0.03cm}
		\tipbox{\scriptsize \textbf{Key Insight:} FATF has issued guidance discouraging indiscriminate de-risking.}
	\end{frame}
	
	\begin{frame}
		\frametitle{\fontsize{11}{13}\selectfont Q78: Alternative Verification Methods}
		\begin{block}{\fontsize{8}{10}\selectfont Topic: Alternative Verification}
			\scriptsize
			What methods can support financial inclusion while meeting AML requirements?
		\end{block}
		\vspace{0.05cm}
		\stepbox{\scriptsize
			\keypoint{Correct Answer:} Alternative identification methods such as biometric verification, digital identity systems, and community-based verification.
		}
		\vspace{0.03cm}
		\tipbox{\scriptsize \textbf{Key Insight:} Technology can bridge the gap between inclusion and compliance.}
	\end{frame}
	
	% ================================================================
	% SLIDE 79-80: SUMMARY AND THANK YOU
	% ================================================================
	\begin{frame}
		\frametitle{\fontsize{11}{13}\selectfont SUMMARY: Key Principles — Domain B}
		\scriptsize
		\begin{block}{\fontsize{9}{11}\selectfont Global AFC Frameworks, Governance, and Regulations:}
			\vspace{0.05cm}
			\scriptsize
			\begin{tabular}{|p{0.35\textwidth}|p{0.60\textwidth}|}
				\hline
				\keypoint{Concept} & \keypoint{Application} \\
				\hline
				FATF Role & Sets global standards; conducts mutual evaluations \\
				\hline
				FSRBs & Regional implementation of FATF standards \\
				\hline
				UN Sanctions & Binding Security Council resolutions (1267, 1373) \\
				\hline
				Regulators, LE, FIUs & Distinct but complementary roles \\
				\hline
				NGO Contributions & CPI, FSI, Basel AML Index, ICAR \\
				\hline
				Regulatory Cooperation & MLATs, Egmont Group, home-host supervision \\
				\hline
				NRAs & Input to EWRA; threat, vulnerability, consequence \\
				\hline
				US Framework & BSA, PATRIOT, FinCEN, OFAC \\
				\hline
				EU Framework & AML Directives, AMLR, AMLA (Frankfurt, 2026) \\
				\hline
				Jurisdictional Awareness & Different requirements; conflicts management \\
				\hline
				PPPs & Timely sharing; joint targeting; disruption \\
				\hline
				Private Sector Collaboration & Industry forums; cross-institution sharing \\
				\hline
				Data Sharing & Legal authority; data protection; tipping-off \\
				\hline
				Data Privacy & GDPR; purpose limitation; cross-border transfer \\
				\hline
				Consumer Protection & Tipping-off; account freezing; fair treatment \\
				\hline
				ESG & Governance; social; environmental crime \\
				\hline
				Financial Inclusion & Proportionality; simplified due diligence \\
				\hline
			\end{tabular}
			\vspace{0.05cm}
		\end{block}
		\vspace{0.05cm}
		\tipbox{\scriptsize \textbf{FINAL LESSON:} Domain B tests understanding of the \textbf{global governance structure} — who sets rules, who enforces them, who implements them, and how they cooperate. Understand the \textit{ecosystem} and the \textit{individual answers} become obvious.}
	\end{frame}
	
	\begin{frame}
		\centering
		\vspace{2.5cm}
		{\fontsize{16}{19}\selectfont \textbf{Thank You}} \\[0.2cm]
		{\fontsize{11}{13}\selectfont Keep learning. Keep growing.} \\[0.15cm]
		\rule{3cm}{0.5pt} \\[0.15cm]
		{\fontsize{8}{10}\selectfont \textit{Understanding global frameworks is the foundation of effective AML/CFT compliance.}}
		\vspace{0.3cm}
		{\fontsize{7}{9}\selectfont \textcolor{darkgray}{End of Domain B Review — 80 Slides}}
	\end{frame}
	
\end{document}